\begin{lemma}[Claim 4.5: the terminal Case 2 coloring]
Let \(G\) be a finite loopless multigraph, let \(e_0\) be the distinguished
uncolored edge, and fix a certified ETT \(\mathcal T\).  Let \(v_n\) be its
terminal supporting vertex.  Fix an exact terminal slot \(n\), an exact
earlier slot \(n^-<n\), and an exact reference slot \(j\), with \(n+1\)
equal to the number of certified rounds.  Let
\(T_j,T_n,T_j(v_n),T'_j\) be tree-sequences, and let
\(\varphi_{j-1},\varphi_n,\sigma_n,\varrho_1\) be the displayed colorings.
Assume that \(\varphi_n,\sigma_n,\varrho_1\) are proper and leave exactly
\(e_0\) uncolored.  The ordered tree data satisfy
\[
 T_j(v_n)\preccurlyeq T_j\preccurlyeq T_n,\qquad
 T_j(v_n)\preccurlyeq T'_j,
\]
the last displayed vertex of \(T_j(v_n)\) is \(v_n\), and \(v_n\) belongs
to the displayed vertex sets required by these inclusions.

Let \(D_{j-1},D,R,S'\) be the preceding control set, the terminal control
set, the restoration-color set, and the support-color set.  Let
\(\gamma_{1_1},\gamma_n,\delta_n\) be colors, let \(f_n\) be an edge, and
let \(x\) be a vertex.  Retain the following Case 2 hypotheses explicitly.
\begin{enumerate}
\item \(T'_j\) is a \(\varrho_1\)-\noderef{TashkinovClosure} of
      \(T_j(v_n)\), \(V(T'_j)\) is elementary, and
      \[
       \overline{\varrho}_1(v)\cap S'=\varnothing
       \quad(v\in V(T'_j)\setminus\{v_n\}),\qquad
       \overline{\varrho}_1(v_n)\cap S'=\{\delta_n\}.          \tag{C2}
      \]
      The colors \(\gamma_n,\delta_n,\gamma_{1_1}\) lie in \(S'\), and
      \(\gamma_n\ne\delta_n\).
\item \(\sigma_n\) is \((T_n,D,\varphi_n)\)-stable,
      \(\gamma_n,\delta_n\in D\), and whenever
      \(\varrho_1(a)\ne\sigma_n(a)\), the exact edge identity \(a\) is not
      internal to \(V(T_j)\) and the old and new colors of \(a\) both lie
      in \(R\).  Neither \(\gamma_n\) nor \(\delta_n\) lies in \(R\).
\item Under \(\varphi_n\), the exact \(\gamma_n\)-boundary of \(T_n\) is
      \(\{f_n\}\); every \(\delta_n\)-boundary edge of \(T_n\) is incident
      with \(T_j(v_n)\); \(f_n\) is incident with \(v_n\); and
      \(\delta_n\in\overline{\varrho}_1(v_n)\).
\item A specified simple \((\gamma_n,\delta_n)\)-path \(P\) from \(v_n\)
      to \(x\notin V(T_n)\cup V(T'_j)\) has support equal to the entire
      two-color component through \(v_n\), and it contains a vertex of
      \(T_n\) other than \(v_n\).
\item The Claim 4.3 data hold in their exact forms: \(\varrho_1\) is
      \((T_j(v_n)-v_n,D_{j-1},\varphi_{j-1})\)-stable;
      \[
       \overline{\varrho}_1(v)=\overline\varphi_{j-1}(v)
       \quad(v\in V(T_j)\setminus\{v_n\}),\qquad
       \overline\varphi_{j-1}(v_n)
       \subseteq\overline{\varrho}_1(v_n)\cup\{\gamma_{1_1}\}; \tag{b2v}
      \]
      \(T_j\) is obtained from \(T_j(v_n)\) by a
      \(\varphi_{j-1}\)-Tashkinov augmentation sequence; every exact edge
      appended after that seed has the same \(\varrho_1\)- and
      \(\varphi_{j-1}\)-color; and
      \(\delta_n\notin\overline\varphi_{j-1}(T_j)\).
\end{enumerate}

Put
\[
 L_1=\noderef{ExtensionBaseSupportClass}(v_n),\qquad
 S'=\bigcup_{t\in L_1}\mathcal T.\mathrm{stored}(t).          \tag{1}
\]
Assume \(n\in L_1\), \(n^-\in L_1\setminus\{n\}\),
\(\delta_{n^-}=\gamma_n\), and that the exact gamma and delta fields at
\(n\) are \(\gamma_n,\delta_n\).  For every exact
\(t\in L_1\setminus\{n\}\), assume that its certified pair
\(\gamma_t,\delta_t\) is \(T_t\)-strongly interchangeable under
\(\varrho_1\), that both colors lie in \(S'\), that neither equals
\(\delta_n\), and that the support-order condition puts every relevant
component vertex other than \(v_n\) inside \(V(T_n)\).

Write the exact increasing enumeration of the whole support class as
\[
 L_1=\{i_0<i_1<\cdots<i_p\},\qquad i_p=n.                    \tag{2}
\]
At each \(i_a\), assume the exact support is \(v_n\), the exact extension
kind is PE, the exact gamma and delta fields are
\(\gamma_a,\delta_a\), and the exact stored set is
\(\{\gamma_a,\delta_a\}\).  Assume
\[
 \gamma_{a+1}=\delta_a\quad(a<p),\qquad
 \gamma_0,\delta_0,\ldots,\delta_p
 \text{ are injectively indexed},\qquad
 \gamma_0=\gamma_{1_1}.                                     \tag{3}
\]
The tree stored at slot \(j\) is exactly \(T_j\), and \(T_j\preccurlyeq
T_t\) for every \(t\in L_1\setminus\{n\}\).

Finally retain, as an explicit ambient Case 2 construction premise, the
present-color fact not contained in the conclusion of Claim 4.6: whenever
\(u\in V(T'_j)\setminus V(T_n)\) and
\(\beta\in\overline{\varrho}_1(u)\), define
\[
 \varrho_0=\noderef{OutsideTreeColorSwap}
   (T'_j,\varrho_1,\beta,\delta_n),\qquad
 \mu_0=\noderef{TwoColorComponentSwap}
   (G,\varrho_0,\gamma_n,\beta,\mathsf{inl}(u)).              \tag{4}
\]
Then every preterminal color \(\gamma_a\), \(a<p\), is present at \(u\)
under \(\mu_0\).  This premise is not attributed to Claim 4.6; it is the
explicit obligation that the later extension-base assembly must discharge
from the preceding Case 2 construction.

Then there are \(u\in V(T'_j)\setminus V(T_n)\) and an exact simple exit
suffix \(Q=P[x,u]\), containing at least one graph-edge identity, such that
\(u\) is the only vertex-node of \(Q\) in \(V(T_n)\cup V(T'_j)\), the
exit has a \(\gamma_n\)-colored boundary identity at \(u\), and \(u\) is
not incident with \(f_n\).  For every
\(\beta\in\overline{\varrho}_1(u)\), there are \(z\in V(T_j)\), at an
exact position \(q_z\ge |T_j(v_n)|\), and a proper coloring
\(\varrho_3\) leaving exactly \(e_0\) uncolored, with the following
properties.  The position \(q_z\) is the first later position missing
\(\beta\), if one exists, and is the terminal position otherwise; moreover
\(\beta\) is missing at no vertex before \(z\) and colors no exact edge
appended between the seed and \(z\).  Finally:
\begin{enumerate}
\item[(d1)] \(\varrho_3\) is
      \((T_j(v_n)-v_n,D_{j-1},\varphi_{j-1})\)-stable.
\item[(d2)] For every
      \(v\in V(T_j(z)\cup T'_j(u))\setminus\{u,z\}\),
      \[
                  \overline{\varrho}_3(v)
                    =\overline{\varrho}_1(v),
      \]
      and every exact displayed edge identity of
      \(T_j(z)\cup T'_j(u)\) has the same color under
      \(\varrho_3\) and \(\varrho_1\).  In addition,
      \[
       \beta\in\overline{\varrho}_1(z)
       \quad\Longrightarrow\quad
       \delta_n\in\overline{\varrho}_3(z).
      \]
\item[(d3)] \(\gamma_{1_1}\in\overline{\varrho}_3(u)\).
\end{enumerate}
All prefix and path assertions above concern the displayed positions and
exact edge identities, so parallel edges are retained distinctly.
\end{lemma}
\begin{proof}
Apply \noderef{ExtensionBaseClaim46} to the stated closure, boundary,
component-path, b1--b4, support-color, and support-order data.  Every
hypothesis is present above individually; in particular, (1) is the exact
support-set identity, and the two exact slots \(n^-,n\) provide the
displayed gamma and delta identities required there.  We obtain a vertex
\(u\in V(T'_j)\setminus V(T_n)\), the asserted exit suffix \(Q\), and all
of its exact path and boundary conclusions.  Fix
\(\beta\in\overline{\varrho}_1(u)\).  The same application gives the
position \(q_z\), the vertex \(z\), its two-case minimality certificate,
and the two exclusions before \(z\).  Define \(\varrho_0,\mu_0\) by (4).
The five conclusions of that node give:
\begin{align}
 &\mu_0\text{ is proper and leaves exactly }e_0\text{ uncolored};
                                                                    \tag{5}\\
 &\mu_0\text{ is }(T_j(v_n)-v_n,D_{j-1},\varphi_{j-1})
    \text{-stable};                                                 \tag{6}\\
 &\overline\mu_0(v)=\overline{\varrho}_1(v)
   \quad(v\in V(T_j(z)\cup T'_j(u))\setminus\{u,z\});               \tag{7}\\
 &\mu_0(f)=\varrho_1(f)
   \quad(f\in E(T_j(z)\cup T'_j(u)));                               \tag{8}\\
 &\beta\in\overline{\varrho}_1(z)
      \Longrightarrow\delta_n\in\overline\mu_0(z),                 \tag{9}\\
 &\gamma_n\in\overline\mu_0(u),\qquad
   \beta\notin\overline\mu_0(u),                                   \tag{10}\\
 &(\gamma_t,\delta_t)\text{ is }T_t\text{-strongly interchangeable
   under }\mu_0\quad(t\in L_1\setminus\{n\}),                       \tag{11}\\
 &\overline\mu_0(T'_j-u)\cap S'
    =\overline\mu_0(v_n)\cap S'=\{\delta_n\},                       \tag{12}\\
 &\mu_0(T'_j.\operatorname{edge}(a))\notin S'\setminus\{\delta_n\}
   \quad(a\ge |T_j(v_n)|).                                         \tag{13}
\end{align}
Equations (5)--(13) are exactly e1--e5 of
\noderef{ExtensionBaseClaim46}; no property from that node's proof is
being used.

We next verify every input of \noderef{ExtensionBaseAlgorithmC}.  The
exact enumeration (2), the slot-support, PE-kind, gamma, delta, stored-pair,
chain, injectivity, and \(T_j\preccurlyeq T_t\) assumptions supply its
finite support-class data.  Equations (5), (10), and (11) supply initial
properness, unique uncoloredness, terminal-gamma missingness, and initial
strong interchangeability.  The explicit final premise of the statement,
not Claim 4.6, supplies the required initial presentness of every
preterminal gamma at \(u\).

It remains only to check Algorithm (C)'s initial missing-color exclusion at
\(v_n\).  Put
\[
 U=\bigcup_{t\in L_1\setminus\{n\}}
       \mathcal T.\mathrm{stored}(t).
\]
By (1), \(U\subseteq S'\).  If
\(c\in\overline\mu_0(v_n)\cap U\), then (12) first gives
\(c=\delta_n\).  Exactness of (2) writes the witnessing slot as \(i_a\).
It cannot be \(i_p=n\), so strict increase gives \(a<p\).  Its exact stored
set is \(\{\gamma_a,\delta_a\}\), whereas the stated earlier-slot
exclusion says that neither member is \(\delta_n\), a contradiction.
Thus
\[
                         \overline\mu_0(v_n)\cap U=\varnothing.       \tag{14}
\]
All inputs of \noderef{ExtensionBaseAlgorithmC} are now verified.  Choose
the resulting \noderef{ExtensionBaseAlgorithmCTrace} \(\mathcal C\),
which processes \(L_1\setminus\{n\}\) in exact decreasing slot order and
satisfies \noderef{ExtensionBaseAlgorithmCInvariant}.  Write its successive
colorings as
\[
 \mu^0=\mu_0,\mu^1,\ldots,\mu^r,\qquad
 \varrho_3=\mu^r.                                                   \tag{15}
\]
By the first clause of the invariant, every \(\mu^a\), in particular
\(\varrho_3\), is proper and leaves exactly \(e_0\) uncolored.  These are
precisely the two coloring-membership assertions in the conclusion; see
\noderef{IsProperEdgeColoring} and \noderef{IsOneEdgeUncolored}.

We prove the preservation needed for (d1) directly at each trace step.
Fix a processed slot \(t\), put
\(\alpha=\gamma_t\), \(\eta=\delta_t\), let \(K\) be the full
\((\alpha,\eta)\)-component rooted at \(u\) before the switch, and let
\(w\in V(T_j(v_n)-v_n)\).  The exact initial-segment field of the invariant
puts \(T_j\preccurlyeq T_t\).  Hence \(w\in V(T_t)\) and \(w\ne v_n\).
The component-intersection clause of
\noderef{ExtensionBaseAlgorithmCInvariant} gives
\(\mathsf{inl}(w)\notin K\).  If an edge incident with \(w\) had its edge
node in \(K\), its pair-color makes it adjacent to \(\mathsf{inl}(w)\)
by \noderef{TwoColorIncidenceGraph}; the connected-component definition
\noderef{TwoColorComponent} would then also put \(\mathsf{inl}(w)\) in
\(K\), a contradiction.  Therefore
\noderef{TwoColorComponentSwap} fixes every exact identity incident with
\(w\).  Expanding \noderef{PresentColors} and \noderef{MissingColors}, it
also fixes the complete missing-color set at \(w\).

Induction over the trace shows that \(\varrho_3\) agrees with \(\mu_0\) on
every edge incident with \(V(T_j(v_n)-v_n)\), and has the same missing set
as \(\mu_0\) at every vertex there.  Thus it is stable on that prefix
relative to \(\mu_0\), directly by
\noderef{StableColoringBeforeVertex}.  Compose these equalities with (6):
an edge selected by the control or missing-color clause for
\(\varphi_{j-1}\) is first fixed by (6) and then by the trace, while each
vertex missing set is equal successively to the \(\mu_0\)- and
\(\varphi_{j-1}\)-missing set.  This is the prefix form of the transitivity
calculation formalized for ordinary tree stability by \noderef{sc2}; the
two clauses have just been checked explicitly because the tablet represents
\(T_j(v_n)-v_n\) by \noderef{StableColoringBeforeVertex}.  This proves
(d1).

We now prove (d2), retaining exact edge identities.  First consider
\(T_j(z)\).  It is an ordered prefix of \(T_j\preccurlyeq T_t\).  If an
exact displayed edge \(f\) of this prefix belonged to \(K\), both endpoint
vertices of \(f\) would also belong to \(K\), by
\noderef{TwoColorIncidenceGraph} and \noderef{TwoColorComponent}.  The
attachment law in \noderef{TreeSequence} puts both endpoints in \(V(T_t)\).
The invariant permits only \(v_n\) from \(V(T_t)\), so both endpoints would
equal \(v_n\), contrary to \noderef{LooplessMultigraph}.  Hence every such
edge identity is fixed by this switch.

For a displayed vertex \(w\in V(T_j(z))\), if \(w\ne v_n\), the same
component-intersection clause shows that no incident edge is switched, so
its missing set is fixed.  If \(w=v_n\) and \(v_n\notin K\), the conclusion
is again immediate.  In the remaining case the invariant says that the
neighbor set of \(v_n\) in the two-color incidence graph is not
subsingleton.  Properness and
\noderef{TwoColorIncidenceDegreeAtMostTwo} bound it by two.  Two distinct
identities of the same color would share \(v_n\), hence be adjacent by
\noderef{AdjacentEdges}, and contradict
\noderef{IsProperEdgeColoring}.  Consequently there is exactly one
incident \(\alpha\)-edge and exactly one incident \(\eta\)-edge, and these
are the two component neighbors.  The switch exchanges their colors and
fixes every other incident identity.  The set of present colors, and
therefore the missing set, is unchanged at \(v_n\).  Induction proves exact
edge-value and vertexwise missing-set preservation on \(T_j(z)\) throughout
the trace.

For \(T'_j(u)\), use (12)--(14) rather than asserting that each component
avoids every prefix vertex.  At a current preterminal slot \(t=i_a\), its
exact stored-pair field and (1) give
\(\alpha,\eta\in S'\), while the earlier-slot exclusion gives
\(\alpha,\eta\ne\delta_n\).  Equation (13) therefore excludes every exact
appended edge of \(T'_j(u)\) from \(K\): its \(\mu_0\)-color is neither
\(\alpha\) nor \(\eta\).  The seed identities are already fixed by the
preceding \(T_j(z)\)-argument, since
\(T_j(v_n)\preccurlyeq T_j\preccurlyeq T_t\).  Thus all exact displayed
edges of \(T'_j(u)\) are fixed at the first step.  The same reasoning is
inductive, since earlier steps have fixed those values.

Let \(w\in V(T'_j(u))\setminus\{u\}\).  Equation (12) says that neither
current switched color is missing at \(w\) under \(\mu_0\).  By induction
the same holds immediately before the current step.  Properness gives one
incident \(\alpha\)-edge and one incident \(\eta\)-edge, each unique by
\noderef{AdjacentEdges} and \noderef{IsProperEdgeColoring}.  If
\(\mathsf{inl}(w)\notin K\), neither is changed.  If
\(\mathsf{inl}(w)\in K\), both incident edge nodes are adjacent to \(w\)
and hence lie in the same full component \(K\); the switch exchanges the
two colors.  In either case the present-color set, and hence the complete
missing-color set, is unchanged at \(w\).  This proves the simultaneous
inductive assertion that the trace fixes every displayed prefix-edge value
and every required prefix-vertex missing set.  Notice that this argument
allows a component to pass through a prefix vertex; it proves the required
nonendpoint behavior from e5 instead of assuming unsupported vertex
avoidance.

Equations (7) and (8) now transfer these final
\(\varrho_3\)-versus-\(\mu_0\) equalities to
\(\varrho_3\)-versus-\(\varrho_1\) equalities on
\(T_j(z)\cup T'_j(u)\), with \(u,z\) excluded exactly as in the statement.
For the conditional clause, (9) makes \(\delta_n\) missing at \(z\) under
\(\mu_0\).  At every trace step the two switched colors belong to an
earlier stored pair and are both different from \(\delta_n\).  Expansion of
\noderef{TwoColorComponentSwap} therefore fixes every
\(\delta_n\)-colored incident value and cannot create one from another
color.  Membership of \(\delta_n\) in the missing set at \(z\) is preserved
through the trace.  This proves all of (d2).

It remains to prove (d3).  The support class contains the strict earlier
slot \(n^-<n\); exactness and strict increase in (2) therefore imply
\(p>0\).  The trace enumerates all preterminal slots and is strictly
decreasing, so its last processed exact slot is \(i_0\), whose gamma is
\(\gamma_0=\gamma_{1_1}\) by (3).  More explicitly, before processing the
largest preterminal position, (10) and the chain equation in (3) say that
its delta is missing at \(u\), while the explicit ambient premise says its
gamma is present.  Properness makes its incident gamma edge unique by
\noderef{AdjacentEdges} and \noderef{IsProperEdgeColoring}.  The rooted
component switch replaces that edge by a delta edge, making that gamma
missing.  Injectivity in (3) says every lower gamma differs from both colors
already switched, so its incident witness remains present; the chain
identity says that the newly missing gamma is precisely the delta required
at the next lower position.  Downward induction repeats this argument at
every exact preterminal slot.  At the final step \(i_0\), it makes
\(\gamma_0=\gamma_{1_1}\) missing at \(u\), and there is no later switch
which could change that conclusion.  By \noderef{MissingColors},
\(\gamma_{1_1}\in\overline{\varrho}_3(u)\), proving (d3) and completing
Claim 4.5.
\end{proof}
